\section*{Discussion}
Our results only partially support our first hypothesis. Covert regulators are able to reduce cheating in protein
design, but barely affect DeepSeek teams working on SWE tasks. Furthermore, in the multi-domain study, a single
explicit rule naming the exploit did at least as well as the presence of any regulator. We were surprised by two
results in particular: a regulator that also works on the task mostly changes its own cheating, and workers took a
teammate's copy of the fix mainly where the regulator had used it.

This suggests regulators need to be more proactive in actively discouraging their peers, and that peer pressure is
likely less impactful when a regulator acts against its own adopted goals. Currently, we measure regulator
effectiveness by its impact on the cheating rate on the task, but not why the cheating rate changed---it could be
the regulator's oversight, or simply the presence of a collaborator. Future work would also analyze whether the
regulator is doing its job by measuring things like coverage (whether the regulator inspects each teammate's
actions) and correction (whether the behavior it flags is actually kept out of the shipped result).

Given that some models have lower propensities to cheat at baseline, they could make for more effective
regulators, even in teams of other models. We're also interested in investigating how regulators constructed via
activation-based methods (e.g., persona selection) could differ from those simply told their instructions in their
system prompt. We have not yet run disregulators, so our second hypothesis remains untested.

We focus on a specific communication topology, with explicit instructions to work collaboratively toward a shared
task, a ready-made communication protocol, and explicit read/write budgets. In reported incidents, they worked as
solo agents and emergently formed a swarm by building their own communication framework. Our swarms are smaller
than those found in the wild, with only 3-4 agents. Some emergent behavior and conformity dynamics may arise only
with large numbers of agents in a swarm.

Our study is currently limited by small sample sizes. Cells hold 3-18 runs, the multi-domain study currently
includes only one task per domain, and the deep dive's regulator shares the workers' model. Our tasks are also
relatively short and were originally designed for single-agent completion. Swarms operating on tasks that require
more complex delegation of work may behave differently. These tasks also include explicit honeypots that may look
like tests, and we do not currently measure evaluation awareness. We do not consider the effects of memory or
compact actions, which occur in real multi-agent systems and can affect information stored/passed through over
time. Extending each of these, in turn, would be informative about how other parameters affect cheating and
peer-pressure propensities.
